\documentclass[12pt, spanish]{article}
\usepackage[spanish]{babel}
\selectlanguage{spanish}
\usepackage{natbib}
\usepackage{url}
\usepackage[utf8x]{inputenc}
\usepackage{graphicx}
\graphicspath{{images/}}
\usepackage{parskip}
\usepackage{fancyhdr}
\usepackage{vmargin}
\usepackage{multirow}
\usepackage{float}
\usepackage{chngpage}

\usepackage{subcaption}

\usepackage{hyperref}
\usepackage[
    type={CC},
    modifier={by-nc-sa},
    version={4.0},
]{doclicense}

\hypersetup{
    colorlinks=true,
    linkcolor=blue,
    filecolor=magenta,      
    urlcolor=cyan,
}

% para codigo
\usepackage{listings}
\usepackage{xcolor}



%% configuración de listings

\definecolor{listing-background}{HTML}{F7F7F7}
\definecolor{listing-rule}{HTML}{B3B2B3}
\definecolor{listing-numbers}{HTML}{B3B2B3}
\definecolor{listing-text-color}{HTML}{000000}
\definecolor{listing-keyword}{HTML}{435489}
\definecolor{listing-identifier}{HTML}{435489}
\definecolor{listing-string}{HTML}{00999A}
\definecolor{listing-comment}{HTML}{8E8E8E}
\definecolor{listing-javadoc-comment}{HTML}{006CA9}

\lstdefinestyle{eisvogel_listing_style}{
  language = C,
%$if(listings-disable-line-numbers)$
%  xleftmargin      = 0.6em,
%  framexleftmargin = 0.4em,
%$else$
  numbers          = left,
  xleftmargin      = 0em,
 framexleftmargin = 0em,
%$endif$
  backgroundcolor  = \color{listing-background},
  basicstyle       = \color{listing-text-color}\small\ttfamily{}\linespread{1.15}, % print whole listing small
  breaklines       = true,
  frame            = single,
  framesep         = 0.19em,
  rulecolor        = \color{listing-rule},
  frameround       = ffff,
  tabsize          = 4,
  numberstyle      = \color{listing-numbers},
  aboveskip        = 1.0em,
  belowskip        = 0.1em,
  abovecaptionskip = 0em,
  belowcaptionskip = 1.0em,
  keywordstyle     = \color{listing-keyword}\bfseries,
  classoffset      = 0,
  sensitive        = true,
  identifierstyle  = \color{listing-identifier},
  commentstyle     = \color{listing-comment},
  morecomment      = [s][\color{listing-javadoc-comment}]{/**}{*/},
  stringstyle      = \color{listing-string},
  showstringspaces = false,
  escapeinside     = {/*@}{@*/}, % Allow LaTeX inside these special comments
  literate         =
  {á}{{\'a}}1 {é}{{\'e}}1 {í}{{\'i}}1 {ó}{{\'o}}1 {ú}{{\'u}}1
  {Á}{{\'A}}1 {É}{{\'E}}1 {Í}{{\'I}}1 {Ó}{{\'O}}1 {Ú}{{\'U}}1
  {à}{{\`a}}1 {è}{{\'e}}1 {ì}{{\`i}}1 {ò}{{\`o}}1 {ù}{{\`u}}1
  {À}{{\`A}}1 {È}{{\'E}}1 {Ì}{{\`I}}1 {Ò}{{\`O}}1 {Ù}{{\`U}}1
  {ä}{{\"a}}1 {ë}{{\"e}}1 {ï}{{\"i}}1 {ö}{{\"o}}1 {ü}{{\"u}}1
  {Ä}{{\"A}}1 {Ë}{{\"E}}1 {Ï}{{\"I}}1 {Ö}{{\"O}}1 {Ü}{{\"U}}1
  {â}{{\^a}}1 {ê}{{\^e}}1 {î}{{\^i}}1 {ô}{{\^o}}1 {û}{{\^u}}1
  {Â}{{\^A}}1 {Ê}{{\^E}}1 {Î}{{\^I}}1 {Ô}{{\^O}}1 {Û}{{\^U}}1
  {œ}{{\oe}}1 {Œ}{{\OE}}1 {æ}{{\ae}}1 {Æ}{{\AE}}1 {ß}{{\ss}}1
  {ç}{{\c c}}1 {Ç}{{\c C}}1 {ø}{{\o}}1 {å}{{\r a}}1 {Å}{{\r A}}1
  {€}{{\EUR}}1 {£}{{\pounds}}1 {«}{{\guillemotleft}}1
  {»}{{\guillemotright}}1 {ñ}{{\~n}}1 {Ñ}{{\~N}}1 {¿}{{?`}}1
  {…}{{\ldots}}1 {≥}{{>=}}1 {≤}{{<=}}1 {„}{{\glqq}}1 {“}{{\grqq}}1
  {”}{{''}}1
}
\lstset{style=eisvogel_listing_style}


\usepackage[default]{sourcesanspro}

\setmarginsrb{2 cm}{1 cm}{2 cm}{2 cm}{1 cm}{1.5 cm}{1 cm}{1.5 cm}

\title{Práctica 2:\\
Algoritmos basados en poblaciones  \hspace{0.05cm} }                           
\author{Antonio David Villegas Yeguas}                             
\date{\today}                                           

\renewcommand*\contentsname{hola}

\makeatletter
\let\thetitle\@title
\let\theauthor\@author
\let\thedate\@date
\makeatother

\pagestyle{fancy}
\fancyhf{}
\rhead{\theauthor}
\lhead{\thetitle}
\cfoot{\thepage}

\begin{document}

%%%%%%%%%%%%%%%%%%%%%%%%%%%%%%%%%%%%%%%%%%%%%%%%%%%%%%%%%%%%%%%%%%%%%%%%%%%%%%%%%%%%%%%%%

\begin{titlepage}
    \centering
    \vspace*{0.3 cm}
    \includegraphics[scale = 0.50]{ugr.png}\\[0.7 cm]
    %\textsc{\LARGE Universidad de Granada}\\[2.0 cm]   
    \textsc{\large 3º CSI 2019/20 - Grupo 1}\\[0.5 cm]            
    \textsc{\large Grado en Ingeniería Informática}\\[0.5 cm]              
    \rule{\linewidth}{0.2 mm} \\[0.2 cm]
    { \huge \bfseries \thetitle}\\
    \rule{\linewidth}{0.2 mm} \\[1 cm]
    
    \begin{minipage}{0.4\textwidth}
        \begin{flushleft} \large
            \emph{Autor:}\\
            \theauthor\\ 
			 \emph{DNI:}\\
            77021623-M
            \end{flushleft}
            \end{minipage}~
            \begin{minipage}{0.4\textwidth}
            \begin{flushright} \large
            \emph{Asignatura: \\
            Metaheurísticas}   \\     
            \emph{Correo:}\\
            advy99@correo.ugr.es           
        \end{flushright}
    \end{minipage}\\[0.5cm]
  
    {\large \thedate}\\[0.5cm]
    {\url{https://github.com/advy99/MH/}}
    {\doclicenseThis}
 	
    \vfill
    
\end{titlepage}

%%%%%%%%%%%%%%%%%%%%%%%%%%%%%%%%%%%%%%%%%%%%%%%%%%%%%%%%%%%%%%%%%%%%%%%%%%%%%%%%%%%%%%%%%

\tableofcontents
\pagebreak

%%%%%%%%%%%%%%%%%%%%%%%%%%%%%%%%%%%%%%%%%%%%%%%%%%%%%%%%%%%%%%%%%%%%%%%%%%%%%%%%%%%%%%%%%


\section{Descripción del problema de la asignación con restricciones.}

El problema de la asignación con restricciones consiste en una generalización del problema de agrupamiento clásico, bastante común en \textit{Machine Learning}.

El problema del agrupamiento clásico es un problema en el que se recibe como entrada las características de un conjunto de elementos y el número de agrupaciones a realizar, para resolver el problema tendremos como objetivo realizar dichas agrupaciones de los distintos elementos con el fin de organizarlos acorde a las características dadas. Llamaremos a estas agrupaciones \textit{Clusters}.

Como extensión a este problema, nosotros trabajaremos sobre el problema de asignación con restricciones (de ahora en adelante PAR). PAR se basa en el problema del agrupamiento clásico, pero añadiendo al problema restricciones entre los propios elementos, es decir, las distintas parejas que podemos formar con los datos tendrán asociadas restricciones de dos tipos:

\begin{itemize}
	\item{Must-Link (ML): Dados dos datos $D_1$ y $D_2$, si estos datos tienen asociada una restricción ML deberán tener asignado el mismo cluster.}
	\item{Cannot-Link (CL): Dados dos datos $D_1$ y $D_2$, si estos datos tienen asociada una restricción CL deberán tener asignado distinto cluster.}
\end{itemize}

Nosotros trataremos estas restricciones como restricciones débiles, es decir, permitiremos que se incumplan pero penalizándolas, por lo tanto, el objetivo para resolver el PAR es minimizar la distancia de los elementos que conforman los distintos clusters, así como asegurarnos que no se incumple ninguna restricción.

Otra restricción fuerte del problema es que todos los clusters tienen que tener al menos un elemento.

Para comprobar que la distancia entre los elementos de los distintos clusters es mínima, tendremos las siguientes características asociadas a un cluster:

\begin{itemize}
	\item{Centroide: Valor promedio de los datos que conforman el cluster. Con este elementos obtendremos la representación del elemento central del cluster.}
	\item{Distancia media intra-cluster: Con este elemento mediremos como de disperso está nuestro cluster, es decir, si los elementos de un cluster están cercanos entre si.}
\end{itemize}

Además también contaremos con distintas características del PAR:

\begin{itemize}
	\item{Desviación general: Será la media de las desviaciones intra-cluster de los distintos clusters que conforman el PAR. Uno de nuestros objetivos será que este valor sea mínimo.}
	\item{\textit{Infeasibility}: Esta característica nos permitirá conocer cuantas restricciones se están incumpliendo en una posible solución del PAR. Otro de nuestros objetivos será que este valor sea mínimo.}
\end{itemize}

\newpage

\section{Descripción de la implementación común y representación del problema para su resolución.}

Para el desarrollo e implementación del programa que resolverá el PAR he decidido usar C++ como lenguaje de programación.

Para la representación del problema he decidido construir dos clases en C++, la clase PAR y la clase Cluster.



\subsection{Clase PAR:}


Esta clase contará con los siguientes atributos:

\begin{itemize}
	\item {Matriz que almacenará valores reales, donde estarán las características de los datos.}
	\item {Vector de objetos tipo Cluster con el que representaremos las agrupaciones a conseguir.}
	\item {Diccionario con las restricciones entre los distintos datos.}
	\item {Desviación general del problema.}
	\item {Mayor distancia entre dos los distintos datos del problema.}
	\item {\textit{Infeasibility} del problema}
\end{itemize}


\subsection{Clase Cluster:}

Con esta clase representaremos los elementos y operaciones internas de un cluster. Tendrá los siguientes atributos:

\begin{itemize}
	\item {Set que almacenará enteros, representando que elementos conforman dicho cluster.}
	\item {Vector de reales que representará el centroide.}
	\item {Valor real que representará la distancia intra-cluster de los elementos que lo conforman.}
	\item {Una referencia a la clase PAR asociada, con la que obtendremos los datos necesarios sin necesidad de duplicarlos.}
\end{itemize}

Es importante mencionar que la declaración de la clase Cluster se realiza dentro de la clase PAR, ya que no tiene sentido crear un cluster sin tener un problema asociado.


\newpage

\subsection{Representación:}

\subsubsection{Datos:}

Con respecto a la representación del problema, los datos son almacenados en una matriz de datos tipo \texttt{double}, cada fila representará un dato, y las columnas representarán las características de dicho dato.

\subsubsection{Restricciones:}

Las restricciones serán almacenadas en un diccionario, donde la clave será la pareja de datos afectada por la restricción y el valor será 1 si la restricción es ML o -1 si es CL. He decidido usar esta representación ya que nos permite almacenar la información de forma eficiente sin tener que almacenar los elementos que no tiene restricciones entre sí, nos permite acceder a los elementos de una forma más rápida que con una lista (aunque no tanto comparado con una matriz, pero gracias al operador \texttt{find} de la clase esto apenas se nota) y podemos recorrer las restricciones secuencialmente de una forma rápida gracias a los iteradores disponibles en la clase diccionario de C++.


Con esta representación de los datos y las restricciones, las distintas operaciones antes comentadas se pueden resolver de forma sencilla generalizando el número de generalizando el número de características, independientemente del problema.


\subsubsection{Solución:}

Para representar la solución he decido, como he comentado anteriormente, que la solución esté compuesta por un vector de objetos tipo cluster, a pesar de que esta representación no será valida para futuras prácticas, es mucho más sencillo y práctico trabajar con ella en la primera práctica, ya que al separar los distintos clusters en su propio objeto podemos realizar las operaciones que afecten a los clusters de forma mucho más rápida y permitiendo la factorización del problema, ya que por ejemplo, seremos capaces de recalcular el centroide de un cluster sin tener que tener en cuenta los demás, o separar los elementos.

La representación dada en clase (vector de enteros de tamaño N siendo N el número de datos, y cada posición del vector se le asocia un número, que será el cluster al que pertenece) es de gran utilidad en la práctica 2, por lo que tenemos una función que nos intercambiará entre estas soluciones, es decir, una función que dado un vector de clusters nos devolverá un vector de enteros con las asignaciones del argumento dado.

\begin{lstlisting}
//vector de ints en el que devolveremos la solución
vector_sol 

Para todo cluster i en el vector de clusters:
	Para todo elemento j en el cluster i:
		vector_sol[j] = i;
		
Devolver vector_sol

\end{lstlisting}


\subsubsection{Operaciones sobre los clusters:}

Los clusters tendrán principalmente dos funciones que podremos realizar:

\begin{itemize}
	\item {Calcular el centroide}
	\item {Calcular la distancia media intra-cluster}
\end{itemize}

Para calcular el centroide basta con recorrer los elementos que conforman ese cluster (disponibles dentro de la clase Cluster) y calcular el punto medio.

\begin{lstlisting}
Para i = 0 hasta el tamaño del centroide (numero de caracteristicas de un dato) :
	centroide[i]=0.0d
	
Para cada elemento i del cluster:
	Para cada carasterictica j del elemento i:
		centroide[j] += problema.datos[i][j];
		
Para cada caracteristica i del centroide:
	centroide[i] /= num_elementos_cluster 
	
	
\end{lstlisting}

Para calcular la distancia media intra-cluster calcularemos el centroide y tras eso la distancia de todos los elementos al centroide.


\begin{lstlisting}
calcular_centroide();
distancia_intra_cluster = 0
	
Para cada elemento i del cluster:
	Para cada carasterictica j del elemento i:
		distancia_intra_cluster += problema.datos[i][j];
		
Para cada caracteristica i del centroide:
	distancia_intra_cluster /= num_elementos_cluster 
	
	
\end{lstlisting}

\newpage

\subsubsection{Operaciones sobre PAR:}

Sobre el problema podremos aplicar las siguientes operaciones:

\begin{itemize}
	\item {Calcular la desviación general.}
	\item {Buscar el cluster en el que se encuentra cierto elemento.}
	\item {Calcular \textit{Infeasibility} del estado actual del PAR.}
	\item {Calcular la distancia entre dos puntos del problema.}
\end{itemize}

Para calcular la desviación general haremos la media de las distintas desviaciones intra-cluster.


\begin{lstlisting}
desviacion_general = 0
	
Para cada elemento i del vector de clusters:
	clusters[i].calcular_desviacion_intra_cluster()
	desviacion_general += cluster[i].desviacion_intra_cluster()
		
desviacion_general /= clusters.size()	
	
\end{lstlisting}

Para buscar un elemento N dado por parámetro en los distintos clusters simplemente recorreremos el vector de clusters usando el operador find de la clase set, al estar ordenados y ser únicos, esto hará que esta operación sea muy rápida.

\begin{lstlisting}
encontrado = false
encontrado_en_cluster = -1

Mientras !encontrado Y para cada elemento i del vector clusters
	Si clusters[i].elementos.find(N) != clusters[i].elementos.end()
		encontrado = true	
		encontrado_en_cluster = i
		
return encontrado_en_cluster
	
\end{lstlisting}


Para calcular  \textit{Infeasibility} del PAR simplemente recorreremos el diccionario de restricciones, si encontramos dos elementos en distinto cluster y son ML sumamos 1 al total, y si encontramos dos elementos en el mismo cluster y tienen la restricción de CL sumamos 1 al total.

Solo comprobaremos si en el diccionario, la pareja de datos el primer elemento es mayor que el segundo. Esto lo hacemos para evitar contar dos veces la misma restricción, por ejemplo, si tenemos la restricción 0 con 1: ML, también tenemos la 1 con 0: ML, así que solo comprobaremos la 1 con 0. A pesar de usar el diccionario he decidido duplicar de esta forma las restricciones ya que a veces necesitaremos acceder sin tener en cuenta el orden de los datos.

\newpage

\begin{lstlisting}
infac = 0

Para cada elemento i de  restricciones
	Si i.first.first > i.first.second
		Si existe una restricción entre i.first.first e i.first.second
			c1 = buscar_elemento(i.first.first)
			c1 = buscar_elemento(i.first.second)
			
			Si c1 == c2 Y i.second == -1
				infac++
			Si c1 != c2 Y i.second == 1
				infac++
		
		
return infac
	
	
\end{lstlisting}




Tanto en PAR como en Cluster tendremos otras operaciones auxiliares relativas a los algoritmos de búsqueda o comparación, que explicaré más adelante.



\subsubsection{Cambio de representación.}

Como hemos comentado, la representación usada es un vector de objetos tipo Cluster, y bajo esa representación hemos resuelto la práctica 1 en la que usamos un algoritmo greedy y una búsqueda local, sin embargo, para trabajar con los algoritmos basados en poblaciones propuestos para la práctica 2 (Algoritmos Genéticos y Meméticos) esta representación no es válida, ya que no disponemos de una forma de cruzar las soluciones de datos, por lo tanto, planteo una segunda representación, en la que una solución será un vector de N posiciones, siendo N el número de elementos del problema, donde cada posición almacena un entero en el intervalo $[0, num_clusters]$ que indicará que cluster tiene asignado el elemento N.


\begin{figure}[H]
  \centering
      $$[0, 1, 2, ... , 2, 1, 1]$$
 		 \caption{Representación de una solución, donde al primer elemento asignamos el cluster 0 y al último el cluster 1}
  		\label{fig:ej_nueva_sol}
\end{figure}

Vemos como esto nos puede traer problemas a la hora de calcular la función objetivo, porque es mucho más costoso actualizar los centroides, calcular la distancia media intra-cluster, etc, sin embargo esto lo solucionaremos de una forma muy simple, tendremos dos funciones, para pasar de una representación a otra, de forma que realizaremos las distintas operaciones en la representación más cómoda para cada operación.

\newpage

\textbf{Paso de vector de clusters a nueva representación:}

\begin{lstlisting}
Devuelve un vector de enteros - recibe un vector de clusters "clusters":
	
	tam_solucion = 0
	
	Para cada cluster i en clusters:
		tam_solucion += i.size()
		
	solucion = vector de enteros (tamaño = tam_solucion)
	
	Para cada cluster i en clusters:
		Para cada elemento j dentro del cluster i:
			solucion[j] = i
			
	devolver solucion
\end{lstlisting}

De esta forma, a cada elemento le asociamos en su correspondiente posición el valor del cluster al que pertenece.

\textbf{Paso de nueva representación a vector de clusters:}

De forma similar podemos hacer la operación contraria:

\begin{lstlisting}
Devuelve un vector de clusters - recibe un vector de enteros solucion_enteros
	solucion = vector de clusters (el númer de clusters lo tenemos almacena en el problema)
	
	Con i desde 0 hasta solucion_enteros.size():
		solucion[solucion_enteros[i]].add_elemento(i)
	
	devolver solucion
\end{lstlisting}

De esta forma, que \texttt{solucion\_enteros[i]} nos dice el cluster donde se encuentra \texttt{i}, luego simplemente lo añadimos a ese cluster en la representación de vector de clusters.



Como hemos comentado, cuando queramos aplicar una operación que sea costosa sobre una representación, pero simple sobre otra, simplemente cambiaremos de una a otra usando estas funciones.

\newpage

\subsection{Calcular \textit{Infeasibility} parcial.}

Para hacer una factorización, además de poder calcular los atributos de los clusters de forma independiente sin necesidad de recalcular la de todos si no son modificados también necesitaremos una función que dado un cluster y un elemento nos calcula cuantas restricciones incumpliría si lo introducimos en dicho cluster.


\begin{lstlisting}
elemento: nos lo dan como argumento
cluster: nos lo dan como argumento

incumplidas = 0


Para todos los elementos i del cluster dado:
	Si i y elemento tienen una restricción Y dicha restricción == -1
		incumplidas++

Para los elementos j de los clusters != cluster:
	Si j y elemento tienen una restricción Y dicha restricción == 1
		incumplidas++



return incumplidas
	
	
\end{lstlisting}


Esta factorización nos servirá para el algoritmo greedy y el algoritmo de búsqueda local, sin embargo no nos servirá con los nuevos algoritmos de la práctica 2, ya que al modificar una solución, esta será totalmente distinta por lo que tendremos que reevaluarla de nuevo.

\newpage

\subsection{Función objetivo.}

En el caso del PAR nuestro objetivo será agrupar los datos en clusters incumpliendo el mínimo de restricciones.

Para cumplir la primera parte intentaremos que los datos estén lo menos dispersos con respecto a su centroide, por lo que intentaremos minimizar el total de las desviaciones intra-cluster, en resumen, \textbf{minimizar la desviación general}. 

Para la segunda parte, intentaremos minimizar el número de restricciones incumplidas penalizando las soluciones que más restricciones incumplan. Esta penalización se basará en sumar a la desviación general un valor entre 0 y la distancia más grande entre los datos del problema.

Esto lo conseguiremos con este factor, al que llamaremos $\lambda$:

$$ \lambda = \frac{D_{max}}{NumRestricciones} $$ 


De forma que si incumplimos todas las restricciones la solución la consideraremos mucho peor que otra con mayor desviación general pero menor restricciones incumplidas.

Nuestra función objetivo será:

$$ f = C + (\textit{Infeasibility} * \lambda) $$ 

Dicho esto podemos implementar la siguiente función:

\begin{lstlisting}
Devolvemos un double:
	devolver desviacion_general + (calcular_infactibilidad() * get_lambda())
\end{lstlisting}

Teniendo ya Lambda calculado al ser un valor que conocemos desde el inicio del problema.

\newpage

\section{Métodos de resolución del problema.}

\subsection{Algoritmo de Búsqueda Local.}

Este algoritmo se basa en partir de una solución inicial aleatoria, a partir de esa solución explorar una solución vecina, y si esta solución vecina es mejor que la actual, moverse a dicha solución.

Este algoritmo no nos asegurará obtener el optimo, solo nos garantiza obtener un mínimo local (puede que este mínimo local sea el optimo, pero no lo podemos asegurar).

El algoritmo solo cambiará una solución por otra mejor, por lo que el punto de inicio será muy importante.

Para desarrollar este algoritmo necesitamos varios componentes:
 
 \subsubsection{Generador de solución inicial aleatoria.}
 
 Primero he desarrollado una función que nos genera una solución aleatoria que cumple con las restricciones fuertes (todos los clusters tienen al menos un elemento).
 
 \begin{lstlisting}
vaciar_clusters()

indices_aleatorios = {0...datos.size()}

indices_aleatorios = random_shuffle(indices)

contador = 0

// nos aseguramos que cada cluster tiene uno al menos
Para cada indice i desde 0 hasta clusters.size():
	clusters[i].insertar_elemento(indices_aleatorios[contador])
	contador++
	
Para contador < indices_aleatorios.size()
	clusters[EnteroAleatorio(0, clusters.size())].insertar_elemento(indices_aleatorios[contador])
	contador++
	

\end{lstlisting}
 
 
\subsubsection{Generador de vecinos}

En nuestro caso, construiremos el vecindario a partir de una solución, modificando un único elemento de cluster, siempre que no se queden clusters vacíos, es decir, si tenemos una solución S, generaremos un vecino $S_1$ eliminando un elemento de un cluster $i$ e insertandolo en otro cluster $j$, siempre que el cluster $i$ no se quede vacío y $i$ sea distinto a $j$.

El elemento a mover y el nuevo cluster lo escogeremos de forma aleatoria, como explicaré más adelante en el desarrollo del algoritmo completo.

 
\subsubsection{Función objetivo} 
 
Una vez tenemos una primera solución, necesitamos conocer como obtener el valor de la función objetivo asociada a esa función.

Cada vez que generemos un vecino recalcularemos su desviación general, por lo que calcular su función objetivo será:

\begin{lstlisting}
funcion_objetivo = get_desviacion_promedio + ( LAMBDA * calcular_infactibilidad() )
 \end{lstlisting}


Sin embargo, para reducir el tiempo de ejecución vamos a factorizar el cálculo de \textit{Infeasibility}, ya que dado un \textit{Infeasibility} de X, el valor para cualquier vecino será:

$$X - \textit{Infeasibility}\_cambios\_salida + \textit{Infeasibility}\_cambios\_entrada$$
 
 Al solo modificar un elemento por vecino, \textit{Infeasibility}\_cambios\_salida  es el número de restricciones que incumplía dicho elemento en el cluster antiguo y \textit{Infeasibility}\_cambios\_entrada es el número de restricciones que incumple el nuevo cluster.
 
 
 
\begin{lstlisting}
infact -= cumple_restricciones(elemento, antiguo_num_cluster)
infact += cumple_restricciones(elemento, nuevo_num_cluster)

funcion_objetivo = get_desviacion_promedio + ( LAMBDA * calcular_infactibilidad() )
 \end{lstlisting}

\subsubsection{Criterio de aceptación.}

Se aceptará un vecino si su función objetivo es menor que la de la solución actual, es decir, seguiremos una estrategia primero el mejor, en cuanto tengamos un mejor candidato nos movemos a el, en lugar de generar todo el vecindario y quedarnos con el mejor del vecindario.

Finalizaremos la búsqueda si no encontramos mejor vecino en todo el vecindario.

\subsubsection{Exploración del entorno.}

La exploración del entorno se realizará de forma aleatoria, cada iteración (no evaluación) reordenaremos de forma aleatoria los indices con los que recorreremos los elementos y los indices con los que recorreremos los clusters. Cada elemento lo probaremos con los distintos clusters (a excepción del propio), por lo que iremos generando los vecinos de cada elemento sin llegar a generar el vecindario completo, solo generamos vecindario hasta que encontramos un mejor vecino. Lo veremos más adelante en el algoritmo.

\subsubsection{Desarrollo del algoritmo.}

El algoritmo de búsqueda local se basará en generar una solución aleatoria, a partir de esa solución aleatoria generar vecinos, evaluarlos hasta que encuentre un mejor vecino, en cuanto encuentre un mejor vecino moverse a este, y parar en caso de que en todo el vecindario no encuentre un mejor vecino o se llegue al límite de evaluaciones, en nuestro caso 100.000.


Con las distintas operaciones podemos codificar el algoritmo de la siguiente forma:

{\small
\begin{lstlisting}
generar_solucion_aleatoria()
calcular_desviacion_general()

LAMBDA = mayor_distancia / restricciones.size()
evaluaciones = 0
encontrado_mejor = false
indices = {0..datos.size()}
indices_clusters = {0..clusters.size()}
infac = calcular_infactibilidad()
infac_vecino = infac
f_objetivo = get_desviacion_general() + (infac * LAMBDA)
f_objetivo_vecino = 0

do:
	encontrado_mejor = false
	reordenar_aleatoriamente(indices)
	reordenar_aleatoriamente(indices_clusters)
	Para cada elemento i de indices Y !encontrado_mejor
		Para cada elemento j del indices_clusters Y !encontrado_mejor
			antiguo_cluster = buscar_elemento(i)
			Si antiguo_cluster != j Y clusters[antiguo_cluster].size() - 1 > 0
				clusters[antiguo_cluster].eliminar_elemento(i)
				infac_vecino -= incumple_restricciones(i, antiguo_cluster)
				infac_vecino += incumple_restricciones(i, j)
				clusters[j].añadir_elemento(i)
				calcular_desviacion_general()
				f_objetivo_vecino = get_desviacion_general() + (infac_vecino * LAMBDA)
				evaluaciones++				
				
				Si f_objetivo_vecino < f_objetivo
					f_objetivo = f_objetivo_vecino
					infac = infac_vecino
					encontrado_mejor = true
				Si NO
					clusters[j].eliminar_elemento(i)
					clusters[antiguo_cluster].añadir_elemento(i)
					infact_vecino = infac

while evaluaciones < TOPE_BL Y encontrado_mejor
 \end{lstlisting}
 }
 
 
Como consideraciones a este algoritmo, vemos como por el operador de vecino es posible que nos estanquemos en mínimos locales de los que no seamos capaces de salir, siendo este el principal problema del algoritmo de búsqueda local.

También mencionar que el escoger la función objetivo puede hacer variar la solución de una forma bastante importante dependiendo del peso que le demos a incumplir las restricciones, como veremos más adelante en los análisis de los resultados y experimentos, ya que de la función objetivo dependerá la penalización por incumplir restricciones por lo que modificando esta función podremos escoger entre mejorar la desviación general o no permitir que se incumplan restricciones.

\newpage

\subsection{Algoritmos evolutivos basados en poblaciones.}

\subsubsection{Introducción a los algoritmos evolutivos.}

Los algoritmos evolutivos basados en poblaciones, como hemos visto en teoría, son algoritmos que se basan en tener un trabajar con un conjunto de soluciones al que llamaremos población, con el objetivo de aplicar distintos operadores para mejorar estas soluciones de la población. A cada una de las soluciones de la población la llamaremos cromosoma, es decir, una población está formada por cromosomas.

De forma similar, como hemos comentado en la nueva representación, la \textbf{representación usada en la práctica 1} (búsqueda local y el algoritmo greedy) \textbf{no nos servirá aquí}, ya que las operaciones que vamos a aplicar a los cromosomas no se pueden aplicar ya que, como explicaremos más adelante, no sería viable obtener la información ni una forma de aplicar estos cambios. Con la nueva representación, usando un vector de enteros en el que el elemento N se representa con la posición N del vector, podemos hablar de genes, que serán los elementos de la solución, los elementos que conforman a un cromosoma.

Como resumen, en un algoritmo evolutivo encontraremos los siguientes elementos:

\begin{enumerate}
	\item \textbf{Gen}: Representación de un elemento de una solución. En la representación un elemento tendrá asociado el cluster al que pertenece.
	\item \textbf{Cromosoma}: Representación de una solución, compuesto por genes.
	\item \textbf{Población}: Representación de un conjunto de cromosomas, es decir, conjunto de soluciones.
\end{enumerate}

El objetivo de esta práctica será ver como con una población, haciendo que los distintos cromosomas interactúen, conseguiremos explorar el espacio de búsqueda, obteniendo soluciones diversas, a la vez que intentaremos que las interacciones entre los cromosomas nos lleven a mejores soluciones intentando seguir el modelo evolutivo, en el que a partir de los cromosomas obtendremos sus hijos, que contendrán información de ambos padres, lo que hará que las nuevas soluciones intenten tomar la mejor parte de ambos a la vez que obtendremos soluciones diversas, evitando los problemas de la búsqueda local, cuyo problema era estancarse en mínimos locales, estos problemas los evitaremos ya que la exploración es mayor como comentaremos más adelante.


\subsubsection{Esquema de los algoritmos evolutivos.}

En esta práctica se nos pide implementar dos tipos de algoritmos:

\begin{itemize}
	\item {Algoritmos Genéticos.}
	\item {Algoritmos Meméticos.}
\end{itemize}

De los cuales vamos a implementar distintas versiones que explicaremos más adelante, sin embargo todos los algoritmos y variaciones seguirán el mismo esquema general de evolución (cada uno tendrá distintas modificaciones, que explicaremos en el desarrollo especifico de cada uno).

\newpage

El esquema evolutivo consta de 4 pasos:

\begin{enumerate}
	\item \textbf{Selección}: A partir de la población actual, seleccionamos a los elementos que formarán parte del proceso evolutivo. Esta selección se suele aplicar usando una ruleta o como en nuestro caso, un torneo binario.
	\item \textbf{Cruce}: Agruparemos los cromosomas de la población por parejas, y aplicando una probabilidad de cruce decidiremos si cruzar las parejas. Existen diversas formas de aplicar el cruce y más adelante explicaremos dos de las que utilizaremos. Un cruce consiste en, dados dos cromosomas (soluciones) crear dos cromosomas hijos, cada uno de estos hijos estará compuesto por una mezcla de sus padres. La idea es dados dos padres que se comportan bien como solución generen hijos que al usar la información del primer padre sumada a la información del segundo, se comporten mejor.
	\item \textbf{Mutación}: Tras aplicar el cruce existirá una pequeña probabilidad de que los genes que conforman los cromosomas muten, es decir, cambien de valor de forma aleatoria. Más adelante explicaremos un método para conseguir esto.
	\item \textbf{Reemplazamiento}: Tras aplicar la selección, el cruce y la mutación, los cromosomas resultantes reemplazarán a la población anterior. Existen varías formas de aplicar el reemplazamiento, ligadas a la forma de realizar la selección y si se decide hacer con elitismo, es decir, mantener el mejor cromosoma de la población en la nueva población.
\end{enumerate}

Para realizar este esquema de evolución necesitamos distintos operadores, que pasaré a explicar en las siguientes secciones.

\subsubsection{Algoritmos evolutivos a implementar.}

En esta práctica vamos a implementar tres algoritmos (sin contar las variaciones de cruce):

\begin{enumerate}
	\item Algoritmos Genéticos Generacionales: Estos algoritmos se basarán totalmente en el esquema de evolución antes comentado, y cada población seleccionada, cruzada y mutada sustituirá por completo a la generación actual.
	\item Algoritmos Genéticos Estacionarios: Estos algoritmos también se basarán totalmente en el esquema de evolución, sin embargo, la población seleccionada será de únicamente dos elementos, por lo que en cada generación, gran parte de la población se mantendrá y solo cambiarán dos cromosomas.
	\item Algoritmos Meméticos: Estos algoritmos no se basan totalmente en el esquema de evolución, lo combinaran con al búsqueda local, aplicándola cada cierto tiempo sobre todos o parte de los cromosomas, de forma que usará la exploración de los algoritmos genéticos para buscar nuevas zonas del espacio con la explotación de la búsqueda local para encontrar los mínimos de dichas zonas del espacio.
\end{enumerate}

Como hemos comentado, los algoritmos genéticos son algoritmos que favorecerán la exploración, por lo que nos permitirán explorar en zonas donde los mínimos sean distintos, pero por el contrario no favorecerán tanto la explotación, por lo que, en teoría, les costará más encontrar el mínimo. El objetivo de está práctica es estudiar si esto ocurre y de que forma controlar el comportamiento de los algoritmos aplicando los distintos operadores.

\subsubsection{Representación para implementar estos algoritmos}

Como hemos comentado en la sección 2, la implementación de la práctica 1 no nos funcionará para aplicar los operadores.

En esta práctica usaremos la nueva representación, pero además, adelantando la condición de parada de los algoritmos, el número de evaluaciones será muy importante, por lo que al almacenar una población almacenaremos con cada cromosoma su valor de la función objetivo, y no tener que perder una evaluación cada vez que queramos conocerlo, haciendo que solo evaluemos cuando sea apliquemos los operadores que modifiquen un cromosoma.

\subsubsection{Operador de selección.}

Este operador se encargará de realizar la selección de la población sobre la que aplicar los cruces y mutaciones. En esta práctica usaremos como operador de selección un torneo binario, es decir, sacaremos dos cromosomas aleatorios, y escogeremos el de mejor función objetivo (en nuestro caso, como queremos minimizar, el de menor función objetivo), cabe destacar que usando este método el cromosoma de peor calidad nunca se escogerá, ya que siempre perderá el torneo, y además es posible que aparezcan más de una copia del mismo cromosoma, haciendo que el mejor cromosoma al ganar siempre el torneo aparecerá con mayor frecuencia, consiguiendo con esto que se seleccione a los mejores cromosomas de la población. Este operador tendrá una pequeña diferencia dependiendo del algoritmo a usar:

\begin{itemize}
	\item Algoritmos Genéticos Generacionales y Algoritmos Meméticos: La población intermedia seleccionada tendrá tantos cromosomas como la población inicial
	\item Algoritmos Genéticos Estacionarios: La población intermedia seleccionada será únicamente de dos elementos.
\end{itemize}

Dicho esto observamos que el operador de selección será común para todos los algoritmos, únicamente variando el número de veces en las que aplicar el torneo binario. Con esto quedaría la siguiente implementación:

\begin{lstlisting}
Devuelve un vector de parejas <cromosoma, valoración> - recibe: vector de parejas <cromosoma, valoracion> pob_inicial, entero con el número de elementos a seleccionar tam

	pob_seleccionada = vector de parejas <cromosoma, valoracion> vacio
	
	Mientras pob_seleccionada.size() < tam:
		primer_candidato = aleatorio (0, pob_ini.size())
		
		// evitamos hacer el torneo entre un elemento y si mismo
		Hacer:
				segundo_candidato =  aleatorio (0, pob_ini.size())
		Mientas segundo_candidato == primer_candidato;
		
		Si pob_ini[primer_candidato].valoracion > poblacion[segunda_candidato].valoracion:
			//buscamos minimizar, si el primer candidato tiene valoracion mayor, el mejor es el segundo
			pob_seleccionada.añadir (pob_ini[segundo_candidato])
		Si no:
			pob_seleccionada.añadir(pob_ini[primer_candidato])


	Devolver pob_seleccionada
\end{lstlisting}

De esta forma conseguiremos seleccionar la población que formará parte del cruce y mutación.


\subsubsection{Operador de reparación.}


\subsubsection{Operador de cruce uniforme.}

Esta será una de las dos formas de aplicar un cruce. Este operador recibirá una población y una probabilidad de cruce, y para evitar tener que generar una gran cantidad de números aleatorios aplicará la esperanza matemática para decidir si una pareja cruza o no. Las parejas de la población, al ya estar seleccionada de forma aleatoria, serán los indices pares con su siguiente indice, es decir, las parejas serán la \{0, 1\}, \{2, 3\}, etc.

 Este operador obtendrá de forma aleatoria N/2 indices del cromosoma aleatorios y sin repetición siendo N el número de genes de los cromosomas, de forma que el hijo generado mantendrá la asignación del primer padre de los elementos obtenidos de forma aleatoria, mientras que la asignación de los otros elementos serán del segundo padre. Como vemos este operador de cruce no da prioridad a ninguno de los dos padres.
 
Este operador recibirá por referencia la población y la probabilidad de cruce, y devolverá el número de evaluaciones realizadas. En este operador usaremos las herramientas explicadas en el cambio de representación, ya que evaluar la solución es mucho más fácil con la representación de la práctica 1.
 
La implementación sería la siguiente:

\begin{lstlisting}
Devuelve el número de evaluaciones - Recibe un vector de parejas <cromosoma, valoracion> pob y una probabilidad de cruce p_cruce

	NUM_PAREJAS = pob.size()/2
	// aplicamos la esperanza matemática
	NUM_CRUCES = NUM_PAREJAS * p_cruce
	
	evaluaciones = 0
	
	indice = 0
	
	// hacemos el doble de cruces ya que cada pareja tiene que generar dos hijos
	Para i = 0 hasta NUM_CRUCES*2; salto para i = +1:
		valores = {}

		// hasta que no generemos la mitad de indices del cromosoma		
		Mientras valores.size() < pob[indice].cromosoma.size()/2:
			valor = aleatorio(0, pob[indice].cromosoma.size())
			
			// si el valor no esta lo añadimos
			Si valores.find(valor ) == valores.end() 
				valores.añadir(valor)
				
		cruce = nuevo vector de tamaño pob[indice].cromosoma.size()
		
		Desde j = 0 hasta pob[indice].cromosoma.size():
			Si j esta en valores:
				cruce[j] = pob[indice].cromosoma[j]
			Si no:
				cruce[j] = pob[indice+1].cromosoma[j]
				
		reparar_cruce(cruce)
		
		Cargar solucion en clusters con solucion_to_clusters(cruce)
		
		hijo = {cruce, funcion_objetivo(cruce)}
		evaluaciones++
		
		pob.añadir(hijo)
		
		// como siempre cruzamos los dos primeros, cuando estamos en i = 1, ya hemos
		// generado y añadido al final los dos hijos de i = 0 e  i = 1, luego borramos los
		// dos primeros elementos de la población
		Si (i % 2 == 1 ):
			pob.erase(pob.begin())
			pob.erase(pob.begin())
					

	Devolver evaluaciones;

\end{lstlisting}

No devolvemos la nueva población ya que al ser una referencia se va modificando mientras se ejecuta el algoritmo.

\subsubsection{Operador de cruce de segmento fijo.}

\subsubsection{Operador de mutación uniforme.}

\subsubsection{Algoritmos Genéticos Generacionales.}

\subsubsection{Algoritmos Genéticos Estacionarios.}

\subsubsection{Algoritmos Meméticos.}

\subsubsection{Implementación.}

\section{Algoritmo de comparación.}

Como algoritmo de comparación usaremos un algoritmo Greedy, basado en una variación del algoritmo k-medias.

\subsection{Algoritmo Greedy.}

El algoritmo se basa en recorrer los indices de forma aleatoria, asignando los elementos al cluster que menos restricciones incumpla, y de entre los que cumplan esta condición al más cercano. Este algoritmo se centrará en tener el menor número de restricciones, aunque la desviación general sea mucho mayor.

\subsubsection{Función para asignar un cluster}

Para resolver el problema de, una vez escogemos un elemento, buscar y asignar un cluster a dicho elemento he diseñado esta función. El principal cometido es calcular el incremento de la \textit{infeasibility} asociada a introducir el elemento en los distintos clusters usando la función para calcular la \textit{infeasibility} parcial antes comentada. Una vez tenemos los clusters que menos restricciones incumplen nos quedamos con el que tenga el centroide a menor distancia y lo asignamos a dicho cluster.


\begin{lstlisting}
elemento = pasado como argumento
distancia = 0
menor_distancia = infinito
menor_restricciones = infinito

// usamos un vector de pares para saber el indice al ordenarlos
Para cada indice i desde 0 hasta clusters.size()
	aumento_infactibilidad.push_back({i, incumple_restricciones(elemento, clusters[i])})

sort(aumento_infactibilidad)
menor_restricciones = aumento_infactibilidad[0].second

Para cada indice i desde 0 hasta aumento_infactibilidad.size() Y aumento_infactibilidad[i].second == menor_restricciones
	distancia = distancia_puntos(clusters[aumento_infactibilidad[i].first.get_centroide(), datos[elemento]])
	Si distancia < menor_distancia
		menor_distancia = distancia
		cluster_menor_distancia = aumento_infactibilidad[i].first


return cluster_menor_distancia
\end{lstlisting}


\subsubsection{Desarrollo del algoritmo.}

Con el uso de esta función y las descritas en la sección común podemos desarrollar el algoritmo de la siguiente forma:

\begin{lstlisting}
limpiamos los clusters
inicializamos los centroides de forma aleatoria

indices = {0 ... datos.size()}

random_shuffle(indices)

hay_cambios = false
cambios = vector de booleanos con tamaño clusters.size() inicializado todo a false

sol_antigua = clusters

do:
	hay_cambios = false
	cambios = {false, ..., false}
	
	Para cada elemento i de indices:
		num_cluster = buscar_cluster(i)
		clusters[num_cluster].add_elemento(i)
		
	Para cada indice i de 0 hasta clusters.size()
		cambios[i] = sol_antigua[i].get_elementos != clusters[i].get_elementos()
		
		Si cambios[i]
			clusters[i].calcular_centroide();
			sol_antigua[i] = clusters[i];
			
		clusters[i].limpiar()
		
	Para cada indice i de 0 hasta cambios.size()
		hay_cambios = hay_cambios || cambios[i]



while hay_cambios

calcular_desviacion_general()
\end{lstlisting}

Notar que está factorizado, de forma que solo se actualicen los centroides que han modificado sus elementos con respecto a la iteración anterior, sin embargo para volver a iterar tenemos que comprobar que al menos un cluster ha cambiado, de ahí actualizar \texttt{hay\_cambios} en función de todos los clusters.

Como hemos comentado, este algoritmo greedy prioriza el minimizar el número de restricciones incumplidas, y entre los valores que menos incumplen, el cluster más cercano, por lo que en un principio podemos pensar que conseguiremos una baja \textit{infeasibility}, sin embargo, como la decisión de introducir un elemento en un cluster influye de cara a los próximos elementos, y al no tener en cuenta próximos elementos ni ser capaces de volver atrás, esto hará que no consigamos esas \textit{infeasibility} tan baja como creíamos.

\section{Proceso de implementación.}

Para la implementación he desarrollado mis propias clases en C++, como adjunto en la carpeta de fuentes. Para las estructuras de datos he utilizado la STL:

\begin{itemize}
	\item Clase map para las restricciones.
	\item Clase vector para los datos y centroides
	\item Clase set para los elementos de los clusters
	\item Clase pair para almacenar las parejas de restricciones, así como para operaciones auxiliares, como el uso en la función auxiliar del algoritmo greedy.
\end{itemize}

Además de estas clases de C++ y las implementadas por mi explicadas a lo largo de este guión también he usado las funciones de generación de números aleatorios dadas por los profesores añadiendo algunas funciones, como por ejemplo generar un número aleatorio entre 0 y un número dado, de forma que el número generado N cumpla que 0 <= N entre otras.

También he utilizado las funciones dadas por los profesores de la asignatura para medir el tiempo.


\subsection{Manual de uso.}

En la carpeta del código fuente existe un fichero Markdown (se puede abrir como texto plano, pero recomiendo un lector de Markdown para facilitar la lectura, editores como Atom o VSCode tienen uno integrado) README.md en el que se explica toda la estructura del código, carpetas, etc.

En esta sección haré un resumen de esto con lo necesario para compilar y ejecutar el programa.

\subsubsection{Compilar:}

Para compilar el programa hay que moverse a la carpeta de fuentes, donde se encuentra el archivo \texttt{Makefile} y ejecutar:

\begin{lstlisting}
make
\end{lstlisting}

Esto nos generará en la carpeta \textit{bin/} el ejecutable.


\subsubsection{Ejecución:}

Podemos lanzar el programa con:

\begin{lstlisting}
./bin/practica1 <fichero_datos> <fichero_restricciones> <num_clusters> <semilla>
\end{lstlisting}

Cada ejecución del programa lanzará todos los algoritmos explicados en esa misma ejecución.

También cabe destacar que en la ruta del fichero de datos se creará un fichero con la extensión \texttt{.out} con la salida y solución del PAR, este fichero contendrá en su nombre la semilla con la que se ha ejecutado, el algoritmo y el conjunto de datos usado.


\newpage

\section{Experimentos y análisis de resultados.}


\subsection{Descripción de los casos.}

En nuestro caso, para el problema del PAR trabajaremos sobre tres conjuntos de datos:

\begin{itemize}
	\item{Iris: Conjunto de datos sobre tres tipos de flores Iris. En este caso tendremos un conjunto de 150 datos y el objetivo será clasificar estas flores según su tipo.}
	\item{Ecoli: Conjunto de datos con características de células, empleadas para predecir la localización de proteínas. En total son 336 datos, de 8 clases distintas.}
	\item{Rand: Conjunto de datos artificial, formado por tres conjuntos de datos bien diferenciados generados a base de distribuciones normales. En total 150 datos con 3 clasificaciones distintas.}
	\item{Newthyroid: Conjunto de datos con medidas cuantitativas tomadas sobre la glándula tiroides de 215 pacientes con 3 clasificaciones distintas.} 
\end{itemize}

Destacar que en nuestro problema las restricciones tendrán un papel muy importante, ya que los conjuntos de datos son muy distintos y estas restricciones son un añadido no original del problema, podría darse el caso de que dos elementos que en principio son de la misma clase estén separados por las restricciones, en caso de tener muy en cuenta las restricciones en nuestra función objetivo, o viceversa, es decir, que realice la asignación de clases sin tener en cuenta las restricciones si a estas no se les da la suficiente importancia.



\subsubsection{Semillas escogidas.}

Las semillas que voy a usar son:

\begin{itemize}
	\item {123452244}
	\item {9398429}
	\item {12321}
	\item {213566}
	\item {3939021}
\end{itemize}

\newpage

\subsection{Resultados obtenidos.}


\newpage


\subsection{Análisis de resultados.	}

\subsection{Otros experimentos realizados.}


\end{document}
